% !TeX program = xelatex
\documentclass[12pt,a4paper,oneside]{report}
\usepackage{iftex}
\ifPDFTeX
  \PackageError{laporan}{Gunakan XeLaTeX atau LuaLaTeX}
    {fontspec tidak mendukung pdfLaTeX.}
\fi
\usepackage{fontspec}
\IfFileExists{fonts/times.ttf}{
  \setmainfont{times.ttf}[Path=fonts/,BoldFont=timesbd.ttf,
    ItalicFont=timesi.ttf,BoldItalicFont=timesbi.ttf]
}{
  \IfFontExistsTF{Times New Roman}{
    \setmainfont{Times New Roman}
  }{
    \setmainfont{TeX Gyre Termes}
  }
}
\usepackage[provide=*,indonesian]{babel}
% Lengkapi identitas kelompok sebelum pengumpulan.
\newcommand{\JenisLaporan}{LAPORAN TUGAS 2}
\newcommand{\KodeMK}{IF4020}
\newcommand{\NamaMK}{Kriptografi}
\newcommand{\JudulLaporan}{Perancangan dan Analisis Brainrot Cipher}
\newcommand{\NamaAlgoritma}{Brainrot Cipher}
\newcommand{\NomorKelompok}{11}
\newcommand{\NamaSatu}{Muhammad Ra'if Alkautsar}
\newcommand{\NIMSatu}{13523011}
\newcommand{\NamaDua}{Muhammad Aditya Rahmadeni}
\newcommand{\NIMDua}{13523028}
\newcommand{\NamaTiga}{Fityatul Haq}
\newcommand{\NIMTiga}{13523116}
\newcommand{\NamaEmpat}{[Nama anggota 4]}
\newcommand{\NIMEmpat}{[NIM 4]}
\newif\ifEmpatAnggota
\EmpatAnggotafalse % Ganti dengan \EmpatAnggotatrue untuk kelompok 4 orang.
\newcommand{\FotoKelompok}{foto-kelompok.png}
\newcommand{\ProgramStudi}{PROGRAM STUDI TEKNIK INFORMATIKA}
\newcommand{\Fakultas}{SEKOLAH TEKNIK ELEKTRO DAN INFORMATIKA}
\newcommand{\Institusi}{INSTITUT TEKNOLOGI BANDUNG}
\newcommand{\AlamatInstitusi}{JL. GANESHA 10, BANDUNG 40132}
\newcommand{\TahunAkademik}{2026/2027}
\newcommand{\Repositori}{https://github.com/mraifalkautsar/Tugas2-IF4020-2026-K11}
% Tautan bonus diisi langsung di bagian/8_lampiran.tex jika tersedia.


% Tata letak dasar dipertahankan dari template laporan sebelumnya.
\usepackage[a4paper,margin=2.54cm,footskip=1.45cm]{geometry}
\usepackage{setspace}
\setstretch{1.15}
\setlength{\parindent}{1.27cm}
\setlength{\parskip}{0.2\baselineskip}
\usepackage{indentfirst,microtype}
% Pertahankan paragraf rata kiri-kanan tanpa memenggal kata.
\hyphenpenalty=10000
\exhyphenpenalty=10000
\tolerance=2000
\emergencystretch=3em
\usepackage{fancyhdr}
\pagestyle{fancy}
\fancyhf{}
\fancyfoot[L]{\small \KodeMK\ \NamaMK}
\fancyfoot[R]{\small \thepage}
\renewcommand{\headrulewidth}{0pt}
\fancypagestyle{plain}{
  \fancyhf{}
  \fancyfoot[L]{\small \KodeMK\ \NamaMK}
  \fancyfoot[R]{\small \thepage}
  \renewcommand{\headrulewidth}{0pt}
}

\usepackage{titlesec,titletoc}
\renewcommand{\thechapter}{\Roman{chapter}}
\renewcommand{\thesection}{\arabic{chapter}.\arabic{section}}
\renewcommand{\thesubsection}{\thesection.\arabic{subsection}}
\renewcommand{\thefigure}{\arabic{chapter}.\arabic{figure}}
\renewcommand{\thetable}{\arabic{chapter}.\arabic{table}}
\renewcommand{\theequation}{\arabic{chapter}.\arabic{equation}}
\titleformat{\chapter}[display]
  {\centering\bfseries\fontsize{16}{20}\selectfont}
  {\chaptername\ \thechapter}{4pt}{}
\titleformat{name=\chapter,numberless}[display]
  {\centering\bfseries\fontsize{16}{20}\selectfont}{}{0pt}{}
\titlespacing*{\chapter}{0pt}{0pt}{18pt}
\titleformat{\section}
  {\bfseries\fontsize{13}{16}\selectfont}{\thesection.}{0.6em}{}
\titleformat{\subsection}
  {\bfseries\normalsize}{\thesubsection.}{0.6em}{}
\titlespacing*{\section}{0pt}{12pt}{5pt}
\titlespacing*{\subsection}{0pt}{8pt}{4pt}
\setcounter{secnumdepth}{2}
\setcounter{tocdepth}{2}
\titlecontents{chapter}[0pt]{\addvspace{6pt}\bfseries}
  {BAB \thecontentslabel\quad}{}
  {\titlerule*[0.6pc]{.}\contentspage}
\titlecontents{section}[2.6em]{}
  {\contentslabel{2.6em}}{}{\titlerule*[0.6pc]{.}\contentspage}
\titlecontents{subsection}[5.8em]{}
  {\contentslabel{3.2em}}{}{\titlerule*[0.6pc]{.}\contentspage}

\usepackage{graphicx,float,booktabs,tabularx,array,longtable}
\graphicspath{{gambar/}}
\newcolumntype{Y}{>{\raggedright\arraybackslash}X}
\usepackage{caption}
\captionsetup{font=small,labelfont=bf,labelsep=period}
\usepackage{amsmath,amssymb,enumitem,xcolor,listings}
\usepackage{tikz}
\usetikzlibrary{arrows.meta,positioning}
\setlist{topsep=4pt,itemsep=2pt,parsep=0pt}
\lstset{basicstyle=\ttfamily\footnotesize,
  numbers=left,numberstyle=\tiny,numbersep=8pt,
  frame=single,rulecolor=\color{black!30},
  backgroundcolor=\color{black!3},breaklines=true,breakatwhitespace=true,
  columns=fullflexible,keepspaces=true,showstringspaces=false,
  tabsize=4,captionpos=b}
\usepackage{url}
\usepackage[hidelinks,unicode,hypertexnames=false]{hyperref}
\urlstyle{same}
\AtBeginDocument{
  \renewcommand{\chaptername}{BAB}
  \renewcommand{\contentsname}{DAFTAR ISI}
  \renewcommand{\listfigurename}{DAFTAR GAMBAR}
  \renewcommand{\listtablename}{DAFTAR TABEL}
  \renewcommand{\bibname}{DAFTAR PUSTAKA}
  \renewcommand{\figurename}{Gambar}
  \renewcommand{\tablename}{Tabel}
  \renewcommand{\lstlistingname}{Kode}
  \renewcommand{\thelstlisting}{\arabic{chapter}.\arabic{lstlisting}}
  \hypersetup{pdftitle={\JudulLaporan},pdfauthor={Kelompok \NomorKelompok}}
}

% Semua bagian yang belum selesai ditandai dengan \Isi{...}.
\newcommand{\Isi}[1]{\textit{[ISI: #1]}}
\newcommand{\GambarAtauPlaceholder}[3]{%
  \IfFileExists{gambar/#1}{%
    \includegraphics[width=#2,height=#3,keepaspectratio]{#1}%
  }{%
    \fbox{\parbox[c][#3][c]{\dimexpr#2-2\fboxsep-2\fboxrule\relax}{%
      \centering\small Tempat gambar\\[6pt]
      \texttt{\detokenize{gambar/}#1}%
    }}%
  }%
}
\newcommand{\BabTanpaNomor}[1]{%
  \chapter*{#1}\phantomsection
  \addcontentsline{toc}{chapter}{#1}%
}

\begin{document}
\begin{titlepage}
  \thispagestyle{empty}\centering
  {\bfseries\fontsize{16}{20}\selectfont \JenisLaporan\ \KodeMK\par}
  \vspace{8pt}
  {\fontsize{16}{20}\selectfont \NamaMK\par}
  \vspace{8pt}
  {\bfseries\fontsize{16}{20}\selectfont \JudulLaporan\par}
  \vspace{6pt}
  {Kelompok \NomorKelompok\par}
  \vspace{0.8cm}
  \GambarAtauPlaceholder{\FotoKelompok}{0.70\textwidth}{5.5cm}\par
  \vspace{0.8cm}
  {\fontsize{14}{18}\selectfont
    Disusun oleh\par\vspace{6pt}
    \begin{tabular}{@{}l@{\hspace{1cm}}l@{}}
      \NamaSatu & \NIMSatu\\
      \NamaDua & \NIMDua\\
      \NamaTiga & \NIMTiga\\
      \ifEmpatAnggota \NamaEmpat & \NIMEmpat\\ \fi
    \end{tabular}\par}
  \vfill
  {\bfseries\fontsize{14}{19}\selectfont
    \ProgramStudi\par \Fakultas\par \Institusi\par
    \AlamatInstitusi\par \TahunAkademik\par}
\end{titlepage}

\pagenumbering{roman}
\begingroup
\setstretch{1.0}\small
\tableofcontents
\clearpage
\listoffigures
\clearpage
\listoftables
\endgroup
\clearpage
\pagenumbering{arabic}
\chapter{PENDAHULUAN}

\section{Latar Belakang}
Tugas 2 IF4020 Kriptografi meminta perancangan \textit{block cipher} baru,
implementasi enkripsi dan dekripsi dalam lima mode operasi, serta pengujian
\textit{avalanche effect}, entropi, dan histogram. Rancangan harus memiliki
substitusi berbasis tabel, permutasi, kunci putaran, dan operasi tambahan
di luar substitusi serta transposisi~\cite{spesifikasi}.

\Isi{Kembangkan motivasi rancangan, masalah yang ingin dijawab, dan alasan
memilih pendekatan Brainrot Cipher.}

\section{Tinjauan Cipher Pembanding}
AES menjadi pembanding yang relevan karena menggunakan blok 128 bit serta
transformasi substitusi dan permutasi~\cite{aes}. Implementasi Brainrot
menggunakan \textit{state} berukuran $4\times4$ byte, sepuluh putaran,
S-box hasil transformasi S-box AES, dan ekspansi kunci dengan pola serupa
AES-128. Parameter permutasi, matriks pencampuran, konstanta penambahan
kunci, dan substitusi dimodifikasi sebagaimana dijelaskan pada Bab III.

\Isi{Uraikan secara singkat persamaan, perbedaan, dan alasan pemilihan
komponen terhadap cipher pembanding. Jelaskan asal ide serta komponen
yang dirancang sendiri sesuai ketentuan tugas; jangan menyimpulkan
keamanan setara AES hanya dari kesamaan struktur.}

\section{Gagasan dan Ruang Lingkup}
\NamaAlgoritma{} menggunakan \textit{Substitution-Permutation Network}
(SPN), blok dan kunci utama sepanjang 128 bit, serta sepuluh putaran.
Program ditulis dalam Go dan menyediakan mode ECB, CBC, CFB, OFB, serta CTR.
Analisis memisahkan pengujian satu blok dari perilaku enkripsi pesan
utuh pada setiap mode.

\Isi{Lengkapi tujuan rancangan, alasan pemilihan parameter, dan batas
ruang lingkup laporan berdasarkan implementasi akhir.}

\chapter{DASAR TEORI}

\section{Block Cipher dan SPN}
\textit{Block cipher} memetakan blok plainteks berukuran tetap menjadi
blok cipherteks dengan kunci rahasia. Untuk setiap kunci $K$, enkripsi
$E_K$ memiliki invers $D_K$ sehingga $D_K(E_K(P))=P$.
SPN mengulang lapisan substitusi, permutasi atau pencampuran, dan
penggabungan kunci. Substitusi nonlinier mendukung \textit{confusion},
sedangkan penyebaran pengaruh byte masukan mendukung \textit{diffusion}.

\Isi{Tambahkan rujukan teori confusion, diffusion, dan SPN yang benar-benar
digunakan. Pertahankan bab ini ringkas.}

\section{S-box, Permutasi, dan Operasi Tambahan}
S-box memetakan nilai byte menggunakan tabel. Pemetaan bijektif memiliki
tabel invers untuk dekripsi. Permutasi mengubah posisi byte tanpa mengubah
nilainya. Operasi XOR, rotasi bit, serta pencampuran pada $\mathrm{GF}(2^8)$
digunakan oleh implementasi Brainrot. Perkalian medan dalam kode memakai
polinomial reduksi $x^8+x^4+x^3+x+1$, yang juga digunakan
AES~\cite{aes}.

\section{Mode Operasi}
ECB, CBC, CFB, OFB, dan CTR adalah lima mode kerahasiaan yang dijelaskan
pada NIST SP 800-38A~\cite{mode}. ECB memproses blok secara independen;
CBC mengaitkan blok melalui cipherteks sebelumnya. CFB memakai umpan balik
cipherteks, OFB memakai umpan balik keluaran cipher, dan CTR mengenkripsi
nilai pencacah. Rincian implementasi dan penanganan blok akhir tersedia
pada Bab IV.

\section{Padding dan Verifikasi Data}
ECB dan CBC memerlukan masukan berupa blok lengkap. Untuk panjang blok
$b$ byte dan pesan $L$ byte, PKCS\#7 menambahkan $p=b-(L\bmod b)$ byte,
masing-masing bernilai $p$. Pesan yang sudah tepat kelipatan blok tetap
mendapat satu blok padding sehingga penghapusannya tidak ambigu~\cite{padding}.
Kesamaan byte hasil dekripsi dengan masukan membuktikan pemulihan data
pada pengujian tersebut. Deteksi perubahan cipherteks oleh pihak lain
memerlukan mekanisme verifikasi tambahan yang harus dijelaskan tersendiri.

\section{Ukuran Statistik}
Jika dua cipherteks $C$ dan $C'$ memiliki $N$ bit, persentase perubahan
bit dinyatakan sebagai
\begin{equation}
  A(C,C')=\frac{\operatorname{HW}(C\oplus C')}{N}\times100\%,
\end{equation}
dengan $\operatorname{HW}$ menyatakan banyaknya bit bernilai satu.
Untuk satu blok Brainrot, $N=128$. Entropi empiris byte dihitung dari
frekuensi relatif $p_v$ setiap nilai byte $v$:
\begin{equation}
  H=-\sum_{v=0}^{255}p_v\log_2 p_v,\qquad 0\log_2 0=0.
\end{equation}
Nilai maksimum distribusi byte adalah 8 bit/byte. Histogram menampilkan
jumlah kemunculan masing-masing nilai $v$. Ketiga ukuran ini memberikan
deskripsi statistik sampel; hasilnya perlu dibahas bersama ukuran sampel
dan perilaku mode operasi.

\chapter{PERANCANGAN BLOCK CIPHER}
\label{bab:perancangan}

% Rincian berikut berasal dari src/cipher/brainrot/, bukan hasil eksperimen.
\section{Parameter dan Representasi State}
\begin{table}[htbp]
  \centering
  \caption{Parameter implementasi Brainrot Cipher.}
  \label{tab:parameter}
  \begin{tabularx}{\linewidth}{@{}lY@{}}
    \toprule Parameter & Nilai\\ \midrule
    Struktur & Substitution-Permutation Network (SPN)\\
    Ukuran blok & 128 bit (16 byte)\\
    Kunci utama & 128 bit (16 byte)\\
    Jumlah putaran & 10\\
    Kunci hasil ekspansi & 11 kunci, $K_0$ sampai $K_{10}$\\
    State & Matriks $4\times4$ byte, diisi per kolom\\
    \bottomrule
  \end{tabularx}
\end{table}
Byte masukan $P_i$ dipetakan ke $s_{r,c}=P_{4c+r}$, untuk $0\leq r,c<4$.
Serialisasi keluaran menggunakan urutan yang sama.

\section{Komponen Transformasi}
\subsection{SubBytes67}
Tabel substitusi berisi 256 nilai dan dibangkitkan melalui
\begin{equation}
  S_{67}(x)=\operatorname{ROTL}_8(S_{\mathrm{AES}}(x\oplus\mathtt{67}),6)
             \oplus\mathtt{43}.
\end{equation}
Konstanta $\mathtt{67}$ dan $\mathtt{43}$ ditulis dalam heksadesimal;
$\mathtt{43}$ sama dengan 67 desimal. S-box AES disimpan di
\path{src/cipher/brainrot/tables.go}. Tabel invers dibentuk dengan
$S_{67}^{-1}[S_{67}[x]]=x$.

\Isi{Jelaskan alasan transformasi S-box dan sifat yang sudah diuji.
Tambahkan tabel S-box lengkap pada lampiran jika diperlukan.}

\subsection{ShiftRows67}
Setiap baris $r$ diputar ke kanan sejauh $o_r=(6+r)\bmod4$ posisi:
\begin{equation}
  s'_{r,(c+o_r)\bmod4}=s_{r,c}.
\end{equation}
\begin{table}[htbp]
  \centering
  \caption{Offset permutasi baris dalam implementasi.}
  \begin{tabular}{@{}lrrrr@{}}
    \toprule Baris $r$ & 0 & 1 & 2 & 3\\
    \midrule Offset kanan $o_r$ & 2 & 3 & 0 & 1\\
    \bottomrule
  \end{tabular}
\end{table}
Inversnya memutar baris ke kiri sejauh offset yang sama.

\subsection{MixColumns67}
Setiap kolom dikalikan dengan matriks $M$ pada $\mathrm{GF}(2^8)$.
Seluruh elemen berikut ditulis dalam heksadesimal:
\begin{equation}
  M=\begin{bmatrix}
    \mathtt{0c}&\mathtt{0a}&\mathtt{06}&\mathtt{06}\\
    \mathtt{07}&\mathtt{0e}&\mathtt{09}&\mathtt{07}\\
    \mathtt{06}&\mathtt{06}&\mathtt{0c}&\mathtt{0a}\\
    \mathtt{09}&\mathtt{07}&\mathtt{07}&\mathtt{0e}
  \end{bmatrix},\quad
  M^{-1}=\begin{bmatrix}
    \mathtt{f4}&\mathtt{ba}&\mathtt{79}&\mathtt{03}\\
    \mathtt{8e}&\mathtt{02}&\mathtt{78}&\mathtt{6a}\\
    \mathtt{79}&\mathtt{03}&\mathtt{f4}&\mathtt{ba}\\
    \mathtt{78}&\mathtt{6a}&\mathtt{8e}&\mathtt{02}
  \end{bmatrix}.
\end{equation}
Penjumlahan medan menggunakan XOR. Perkalian menggunakan pergeseran bit
dan reduksi dengan konstanta $\mathtt{1b}$ setelah bit tertinggi melimpah.

\Isi{Jelaskan pemilihan matriks dan bukti invers atau sifat difusi yang
telah diperiksa.}

\subsection{AddRoundKey67}
Untuk kunci putaran $K_j$, setiap elemen diperbarui sebagai
\begin{equation}
  s'_{r,c}=s_{r,c}\oplus (K_j)_{4c+r}\oplus\mathtt{43}.
\end{equation}
Transformasi ini merupakan invers dirinya sendiri.

\section{Alur Enkripsi dan Dekripsi}
\begin{figure}[htbp]
  \centering
  \begin{tikzpicture}[
    node distance=5mm,
    box/.style={draw,rounded corners,align=center,text width=9cm,
      minimum height=8mm,font=\small},
    flow/.style={-{Latex},thick}]
    \node[box] (in) {Plainteks 16 byte $\rightarrow$ state $4\times4$};
    \node[box,below=of in] (initial) {AddRoundKey67 dengan $K_0$};
    \node[box,below=of initial] (rounds) {Putaran $j=1,\ldots,9$\\
      SubBytes67 $\rightarrow$ ShiftRows67 $\rightarrow$ MixColumns67\\
      $\rightarrow$ AddRoundKey67 dengan $K_j$};
    \node[box,below=of rounds] (final) {Putaran 10\\
      SubBytes67 $\rightarrow$ ShiftRows67 $\rightarrow$ AddRoundKey67
      dengan $K_{10}$};
    \node[box,below=of final] (out) {Serialisasi state $\rightarrow$ cipherteks 16 byte};
    \draw[flow] (in) -- (initial);
    \draw[flow] (initial) -- (rounds);
    \draw[flow] (rounds) -- (final);
    \draw[flow] (final) -- (out);
  \end{tikzpicture}
  \caption{Alur enkripsi satu blok sesuai implementasi saat ini.}
  \label{fig:alur-enkripsi}
\end{figure}
Enkripsi dimulai dengan AddRoundKey67 menggunakan $K_0$.
Putaran 1--9 menjalankan keempat transformasi; putaran terakhir melewati
MixColumns67, seperti pada Gambar~\ref{fig:alur-enkripsi}.

Dekripsi menjalankan AddRoundKey67 dengan $K_{10}$, lalu
InverseShiftRows67 dan InverseSubBytes67. Untuk $j=9,\ldots,1$, urutannya
adalah AddRoundKey67 dengan $K_j$, InverseMixColumns67,
InverseShiftRows67, dan InverseSubBytes67. Proses ditutup dengan
AddRoundKey67 menggunakan $K_0$.

\Isi{Tambahkan contoh state dan kunci untuk satu blok, serta justifikasi
jumlah putaran dan peniadaan pencampuran pada putaran terakhir.}

\section{Pembentukan Kunci Putaran}
Kunci utama dibagi menjadi empat word $w_0,\ldots,w_3$, masing-masing
empat byte. Ekspansi menghasilkan 44 word. Untuk $i=4,\ldots,43$,
mula-mula $t=w_{i-1}$. Jika $i\bmod4=0$, $t$ dirotasi satu byte ke kiri,
setiap byte disubstitusi dengan $S_{67}$, lalu byte pertama di-XOR dengan
$\mathrm{Rcon}_{i/4-1}\oplus\mathtt{67}$. Selanjutnya,
\begin{equation}
  w_i=w_{i-4}\oplus t,\qquad
  K_j=w_{4j}\Vert w_{4j+1}\Vert w_{4j+2}\Vert w_{4j+3}.
\end{equation}
Urutan Rcon adalah $\mathtt{01,02,04,08,10,20,40,80,1b,36}$.
Simbol $\Vert$ menyatakan konkatenasi byte; $K_0$ adalah kunci utama.

\Isi{Tambahkan contoh pembentukan kunci dan alasan pemilihan konstanta.
Jelaskan kontribusi operasi tambahan terhadap rancangan.}

\chapter{IMPLEMENTASI DAN SIMULASI}

\section{Struktur Program dan Dependensi}
Program menggunakan Go 1.23 atau lebih baru. Cipher dan mode operasi
diimplementasikan dalam repositori tanpa pustaka kriptografi pihak ketiga.
Pustaka standar \texttt{crypto/rand} dipakai oleh CLI untuk membangkitkan
IV; \texttt{math/rand} hanya dipakai untuk eksperimen yang dapat diulang.

\begin{table}[htbp]
  \centering\small
  \caption{Komponen kode sumber.}
  \begin{tabularx}{\linewidth}{@{}p{0.48\linewidth}Y@{}}
    \toprule Lokasi & Peran\\ \midrule
    \path{src/cipher/cipher.go} & Kontrak Block, Factory, dan Tracer.\\
    \path{src/cipher/brainrot/} & Cipher, transformasi, tabel, dan ekspansi kunci.\\
    \path{src/modes/modes.go} & Lima mode, PKCS\#7, dan pencacah CTR.\\
    \path{src/cmd/crypt/main.go} & CLI enkripsi teks dan dekripsi hex.\\
    \path{src/cmd/analyze/main.go} & CLI analisis berkas uji dan keluaran JSON.\\
    \path{src/analysis/} & Statistik distribusi dan eksperimen perubahan bit.\\
    \path{src/tests/} & Pengujian infrastruktur dengan cipher XOR uji.\\
    \bottomrule
  \end{tabularx}
\end{table}

\section{Antarmuka dan Format Data}
CLI \texttt{crypt} menerima \texttt{-op}, \texttt{-mode},
\texttt{-key-hex}, dan \texttt{-input}. Kunci harus berupa 32 karakter
heksadesimal. Saat enkripsi, masukan berupa teks; saat dekripsi,
masukan berupa hex hasil enkripsi. Mode selain ECB menyimpan IV 16 byte
di depan badan cipherteks: $\mathrm{IV}\Vert C$. Mode ECB hanya menyimpan $C$.
Mode dan kunci tidak disertakan dalam keluaran; pengguna memasukkan
keduanya lagi saat dekripsi.

\begin{lstlisting}[language={},numbers=none,caption={Contoh simulasi CBC dari root repositori pada Bash.}]
key_hex=00112233445566778899aabbccddeeff
ciphertext=$(go -C src run ./cmd/crypt -op encrypt \
  -mode CBC -key-hex "$key_hex" -input "Halo Brainrot")
go -C src run ./cmd/crypt -op decrypt \
  -mode CBC -key-hex "$key_hex" -input "$ciphertext"
\end{lstlisting}

\Isi{Cantumkan keluaran aktual atau tangkapan terminal. Ulangi simulasi
untuk ECB, CBC, CFB, OFB, dan CTR. Deskripsikan antarmuka akhir jika berubah.}

\section{Mode Operasi, IV, dan Counter}
\begin{table}[htbp]
  \centering\small
  \caption{Perilaku mode operasi dalam kode saat ini.}
  \label{tab:mode}
  \begin{tabularx}{\linewidth}{@{}lp{0.23\linewidth}Y@{}}
    \toprule Mode & IV / counter & Blok akhir dan umpan balik\\ \midrule
    ECB & Tidak digunakan & Blok independen; PKCS\#7.\\
    CBC & IV 16 byte & XOR dengan blok cipherteks sebelumnya; PKCS\#7.\\
    CFB & IV 16 byte & Umpan balik blok penuh (CFB-128); potongan akhir
      memakai sebagian keystream tanpa padding.\\
    OFB & IV 16 byte & Umpan balik keluaran cipher; tanpa padding.\\
    CTR & Counter awal 16 byte & Pencacah big-endian 128 bit;
      menolak overflow jika masih ada data; tanpa padding.\\
    \bottomrule
  \end{tabularx}
\end{table}
Untuk CBC, CFB, dan OFB, IV dibuat oleh CLI menggunakan
\texttt{crypto/rand}. Nilai yang sama menjadi counter awal CTR.
Dekripsi CFB, OFB, dan CTR menggunakan fungsi enkripsi blok untuk
membangkitkan keystream.

\Isi{Jelaskan alasan pengelolaan IV/counter dan mekanisme penyimpanannya
pada format berkas akhir. Catat cara mencegah penggunaan ulang keystream.}

\section{Penanganan Padding}
Untuk ECB dan CBC dengan $b=16$, panjang badan cipherteks adalah
$16(\lfloor L/16\rfloor+1)$ byte. Untuk CFB, OFB, dan CTR, panjang badan
cipherteks sama dengan panjang plainteks. Panjang IV tidak termasuk dalam
kedua ukuran tersebut. Dekripsi ECB/CBC memeriksa panjang cipherteks,
nilai byte padding terakhir, dan kesamaan seluruh byte padding.

\Isi{Justifikasi PKCS\#7 dan lengkapi bukti untuk masukan 0, 1, 15, 16,
17, dan 32 byte. Catat panjang plainteks, padding, badan cipherteks,
dan hasil pemulihan; sertakan kasus padding tidak valid.}

\section{Simulasi Berkas Teks dan Biner}
Fungsi mode menerima slice byte, sedangkan CLI analisis membaca berkas.
CLI \texttt{crypt} saat ini hanya menerima teks dari argumen dan belum
menyediakan pembacaan serta penulisan berkas biner sembarang.

\Isi{Lengkapi bagian ini setelah alur enkripsi/dekripsi berkas tersedia.
Gunakan berkas teks dan biner pada seluruh mode; jangan menganggap
simulasi argumen teks sebagai bukti dukungan berkas biner.}

\begin{table}[htbp]
  \centering\small
  \caption{Isian bukti pemulihan data untuk seluruh mode.}
  \label{tab:simulasi}
  \begin{tabularx}{\linewidth}{@{}lYYYY@{}}
    \toprule Mode & Teks: berkas / ukuran & Biner: berkas / ukuran &
      Panjang keluaran & Byte hasil dekripsi sama\\ \midrule
    ECB & --- & --- & --- & ---\\
    CBC & --- & --- & --- & ---\\
    CFB & --- & --- & --- & ---\\
    OFB & --- & --- & --- & ---\\
    CTR & --- & --- & --- & ---\\
    \bottomrule
  \end{tabularx}
\end{table}

\section{Verifikasi Integritas dan Kasus Gagal}
Analisis menggunakan \texttt{bytes.Equal} untuk memeriksa bahwa hasil
dekripsi sama persis dengan masukan. Kode saat ini belum menyimpan
tag autentikasi atau mekanisme verifikasi perubahan cipherteks. Pemeriksaan
padding hanya memvalidasi bentuk padding, sehingga tidak mencakup
seluruh perubahan pada isi pesan.

\Isi{Jelaskan mekanisme integritas yang dipilih pada implementasi akhir
sesuai spesifikasi. Bedakan bukti round-trip dengan deteksi modifikasi
data. Sertakan hasil pengujian kunci salah, data berubah, IV berubah,
data terpotong, dan masukan tidak valid.}

\chapter{ANALISIS KEAMANAN}
% Generated by scripts/render_brainrot_results.py from recorded v2 JSON.
\newcommand{\VTwoBlockRows}{%
1 bit plainteks & 256 & 49.643 & 47 & 79\\
1 bit kunci & 256 & 50.177 & 49 & 82\\
}

% Generated by scripts/render_brainrot_results.py from recorded v2 JSON.
\newcommand{\VTwoModeRows}{%
ECB & 416 & 1.931 & 50.001\\
CBC & 416 & 25.969 & 50.056\\
CFB & 404 & 24.204 & 50.000\\
OFB & 404 & 0.031 & 49.923\\
CTR & 404 & 0.031 & 49.959\\
}

% Generated by scripts/render_brainrot_results.py from recorded v2 JSON.
\newcommand{\VTwoEntropyRows}{%
ECB & 416 & 4.860 & 7.217\\
CBC & 416 & 4.860 & 7.511\\
CFB & 404 & 4.860 & 7.456\\
OFB & 404 & 4.860 & 7.463\\
CTR & 404 & 4.860 & 7.513\\
}

% Generated by scripts/render_brainrot_results.py from recorded v2 JSON.
\newcommand{\VTwoBinaryEntropyRows}{%
ECB & 1040 & 8.000 & 7.197\\
CBC & 1040 & 8.000 & 7.796\\
CFB & 1027 & 8.000 & 7.777\\
OFB & 1027 & 8.000 & 7.785\\
CTR & 1027 & 8.000 & 7.807\\
}


\section{Konfigurasi dan Reproduksibilitas}
CLI \texttt{analyze} menggunakan Brainrot v2 secara default. Masukan bawaan
adalah \path{src/testdata/fixtures/one-block.txt} sepanjang 16 byte dan
\path{src/testdata/fixtures/repeated-text.txt} sepanjang 404 byte.
Parameter bawaan adalah 256 percobaan, seed 42, serta kunci
\texttt{00112233445566778899aabbccddeeff}. Seed hanya digunakan untuk
eksperimen; IV pada CLI enkripsi dibangkitkan secara terpisah.

\begin{lstlisting}[language={},numbers=none,caption={Menjalankan analisis dari root repositori.}]
go -C src run ./cmd/analyze -cipher brainrot -modes all \
  -trials 256 -seed 42 -output ../reports/brainrot-v2/text-seed42.json
go -C src test ./...
\end{lstlisting}

Laporan JSON menyimpan ukuran blok dan kunci, seed, jumlah percobaan,
avalanche satu blok, hasil round-trip per mode, entropi, histogram 256 bin,
dan avalanche pesan utuh. Pengukuran mode hanya memakai badan cipherteks;
IV tidak masuk histogram atau penyebut persentase perubahan bit.

Hasil yang ditampilkan berasal dari \path{reports/brainrot-v2/text-seed42.json}
dan \path{reports/brainrot-v2/binary-seed42.json}. Fixture biner
\path{src/testdata/fixtures/byte-pattern.bin} terdiri atas nilai byte 0--255
yang diulang empat kali, diikuti $\mathtt{00,ff,80}$; totalnya 1027 byte.
Manifest pada folder hasil mencatat kunci, seed, ukuran dan SHA-256 fixture,
versi Go, sistem operasi, commit dasar, serta hash sumber v2. Dua eksperimen
tambahan memakai seed 1337 dengan kunci nol pada teks dan kunci semua satu
pada biner. Seluruh angka adalah pengamatan sampel, bukan ambang kelulusan
keamanan.

\section{Avalanche Effect Satu Blok}
Pada setiap percobaan, satu bit plainteks diubah dengan kunci tetap;
secara terpisah, satu bit kunci diubah dengan plainteks tetap. Percobaan
pertama memakai fixture satu blok, sedangkan berikutnya memakai blok
pseudorandom dari seed yang tercatat. Posisi bit digilir melalui seluruh
128 bit masukan; setidaknya 128 percobaan diperlukan untuk mencakup
setiap posisi. Keluaran mencatat rata-rata, minimum, maksimum,
serta rata-rata dan jumlah percobaan per posisi bit.

\begin{table}[htbp]
  \centering\small
  \caption{Hasil avalanche Brainrot v2 satu blok, 256 percobaan dan seed 42.}
  \begin{tabularx}{\linewidth}{@{}lYYYY@{}}
    \toprule Perubahan & Percobaan & Rata-rata (\%) & Minimum (bit) & Maksimum (bit)\\
    \midrule
    \VTwoBlockRows
    \bottomrule
  \end{tabularx}
\end{table}

Dengan 256 percobaan, setiap posisi bit diuji dua kali. Rata-rata perubahan
bit plainteks dan kunci mendekati separuh keluaran pada sampel ini.
\Isi{Kembangkan pembahasan variasi per posisi bit dari JSON dan
bandingkan dengan eksperimen kunci/seed lainnya.}

\section{Avalanche Effect pada Seluruh Mode}
Untuk pesan utuh, posisi bit plainteks dipilih pseudorandom dan posisi
bit kunci digilir. Pasangan enkripsi sebelum dan sesudah perubahan
menggunakan IV yang sama agar perbedaan berasal dari bit yang diubah.
Penyebutnya adalah jumlah bit seluruh badan cipherteks, termasuk padding
ECB/CBC, sehingga berbeda dari pengujian satu blok.

\begin{table}[htbp]
  \centering\small
  \caption{Avalanche pesan teks 404 byte, 256 percobaan dan seed 42.}
  \begin{tabularx}{\linewidth}{@{}lYYY@{}}
    \toprule Mode & Panjang badan (byte) & Perubahan plainteks (\%) & Perubahan kunci (\%)\\
    \midrule
    \VTwoModeRows
    \bottomrule
  \end{tabularx}
\end{table}
Pada OFB dan CTR dengan kunci serta IV tetap, satu bit plainteks yang
diubah menghasilkan satu bit cipherteks yang berubah karena operasi XOR
dengan keystream yang sama. Pada ECB perubahan terbatas pada blok terkait;
pada CBC dan CFB pengaruhnya bergantung pada posisi perubahan dan umpan
balik. Interpretasi hasil harus membedakan sifat mode dari difusi cipher
satu blok.

\Isi{Bahas hasil kelima mode, posisi bit yang dipilih, pengaruh panjang
pesan, dan perbedaan perubahan plainteks terhadap perubahan kunci.}

\section{Analisis Entropi}
\begin{table}[htbp]
  \centering\small
  \caption{Entropi empiris untuk fixture teks, seed 42; ukuran adalah badan cipherteks.}
  \begin{tabularx}{\linewidth}{@{}lYYY@{}}
    \toprule Mode & Ukuran sampel (byte) & Plainteks (bit/byte) & Cipherteks (bit/byte)\\
    \midrule
    \VTwoEntropyRows
    \bottomrule
  \end{tabularx}
\end{table}
\begin{table}[htbp]
  \centering\small
  \caption{Entropi empiris untuk fixture biner, seed 42.}
  \begin{tabularx}{\linewidth}{@{}lYYY@{}}
    \toprule Mode & Ukuran badan (byte) & Plainteks (bit/byte) & Cipherteks (bit/byte)\\
    \midrule
    \VTwoBinaryEntropyRows
    \bottomrule
  \end{tabularx}
\end{table}
\Isi{Jelaskan pengaruh ukuran sampel, pola plainteks, dan padding.
Nilai mendekati 8 bit/byte saja tidak membuktikan keamanan kriptografis.}

\section{Analisis Histogram}
\begin{figure}[htbp]
  \centering
  \GambarAtauPlaceholder{histogram-semua-mode.png}{\linewidth}{11cm}
  \caption{Histogram fixture teks dan cipherteks lima mode pada seed 42;
    sumbu y memakai skala berbeda dan garis putus-putus menandai distribusi seragam.}
  \label{fig:histogram}
\end{figure}
\begin{figure}[htbp]
  \centering
  \GambarAtauPlaceholder{histogram-biner.png}{\linewidth}{11cm}
  \caption{Histogram fixture biner dan cipherteks lima mode, seed 42;
    sumbu y memakai skala berbeda.}
\end{figure}
Grafik memakai proporsi byte agar panjang badan yang berbeda dapat
dibandingkan. Skala y berbeda antar panel dan tercantum pada setiap sumbu.
\Isi{Bahas pola berulang, khususnya ECB, serta kedekatan distribusi empiris
terhadap distribusi seragam dengan memperhatikan ukuran sampel.}

\section{Pengujian Tambahan dan Keterbatasan}
\Isi{Jika dilakukan, tambahkan pengujian ukuran kecil/sedang/besar,
waktu proses, penggunaan memori, perubahan per putaran melalui Tracer,
atau kajian struktur S-box dan matriks. Sertakan metode serta hasil aktual.}

Pengujian otomatis mempertahankan \texttt{TEST-ONLY-XOR} untuk hasil
chaining yang dapat diprediksi, dan menambahkan pengujian khusus Brainrot
v2. Empat vektor referensi independen memeriksa cipherteks, 21 subkunci
per kasus, serta 66 snapshot. Invers, input berpola dan pseudorandom,
overlap buffer, pemetaan semua 128 bit BPT, serta seluruh 65.536 pasangan
byte/subkunci untuk penjumlahan diperiksa. Lima mode juga diuji dengan
Brainrot pada masukan kosong, batas blok, padding tidak valid, dan overflow CTR.

Pemeriksaan rank biner menguji 36.416 pasangan support untuk setiap
jenis propagasi dan membuktikan branch number byte minimal 5 pada
lapisan ShiftRows67--MixColumns67--BPT revisi untuk perbedaan dan masker
linier. Referensi Python memeriksa batas yang sama secara independen.
Batas tersebut tidak mencakup FCM, penjumlahan, atau seluruh cipher.

Benchmark tiga kali pada Linux amd64, AMD Ryzen 5 5600H, menghasilkan
median enkripsi 20.225 ns/blok pada cipher lama dan 24.959 ns/blok pada
v2 (kenaikan waktu sekitar 23,4\%). Median dekripsi berubah dari
20.141 menjadi 24.673 ns/blok (sekitar 22,5\%). Tidak ada alokasi heap
per blok pada kedua versi. Angka tersebut berasal dari benchmark lokal
dan bukan batas performa universal; keluaran mentah disimpan bersama hasil.

FCM memiliki simetri rotasi pada kedua half, dan fungsi internalnya
memiliki empat kemungkinan perbedaan keluaran dengan probabilitas
$1/4$ untuk perbedaan masukan satu bit pada input seragam. Penjumlahan
byte mempertahankan perbedaan XOR bit tertinggi untuk setiap subkunci.
Pasangan subkunci penjumlahan/XOR $(a,b)$ dan
$((a+128)\bmod256,b\oplus128)$ juga menghasilkan pemetaan byte sama.
Ini adalah sifat komponen, bukan bukti serangan pada cipher lengkap
atau adanya kunci utama ekuivalen.

\Isi{Nyatakan cakupan pengujian akhir, keterbatasan sampel, dan analisis
yang belum dilakukan. Kaitkan kesimpulan hanya dengan bukti yang tersedia.}

\chapter{KESIMPULAN DAN SARAN PENGEMBANGAN}

\section{Kesimpulan}
Berdasarkan perancangan, implementasi, serta analisis yang telah dilakukan terhadap Brainrot Cipher, ditarik beberapa kesimpulan sebagai berikut:
\begin{enumerate}
  \item \textbf{Pemenuhan Spesifikasi Perancangan Cipher:} Algoritma Brainrot berhasil dirancang dan diimplementasikan sebagai \textit{block cipher} simetris 128 bit dengan panjang kunci 128 bit, 10 putaran, serta 21 subkunci terikat melalui prosedur \textit{key schedule}. Struktur berbasis \textit{Substitution--Permutation Network} (SPN) memadukan substitusi nonlinier Kotak-S $S_{67}$ (permutasi matriks $16\times16$ sepanjang 67 iterasi), pergeseran baris ShiftRows67, difusi kolom linear MixColumns67 menggunakan matriks MDS pada $\mathrm{GF}(2^8)$ dengan \textit{branch number} 5, Bit-Plane Transpose (BPT), langkah nonlinier Feistel Cross-Mixing (FCM) pada putaran genap, serta kombinasi subkunci penjumlahan modulo byte dan XOR AddRoundKey67. Seluruh komponen transformasi terbukti dapat dibalik (\textit{invertible}) dan kesesuaian alurnya telah divalidasi menggunakan vektor uji independen serta pencocokan 66 \textit{state} perantara.

  \item \textbf{Implementasi Mode Operasi dan Penanganan Padding:} Kelima mode operasi (ECB, CBC, CFB, OFB, dan CTR) berhasil diimplementasikan untuk memproses data teks (404 byte) dan berkas biner (1027 byte). Skema padding PKCS\#7 pada mode ECB dan CBC berhasil menangani sisa byte maupun pesan yang sudah tepat kelipatan blok, serta menolak padding tidak valid dengan respons HTTP 400. Mode CFB, OFB, dan CTR memproses sisa byte melalui operasi XOR dengan potongan \textit{keystream} tanpa memerlukan padding. Seluruh sepuluh kombinasi pengujian \textit{round-trip} menunjukkan pemulihan data yang identik dengan masukan aslinya melalui kesamaan byte langsung dan nilai \textit{digest} SHA-256.

  \item \textbf{Batas Verifikasi dan Autentikasi Integritas:} Keberhasilan pemulihan data pada pengujian \textit{round-trip} membuktikan invertibilitas algoritma dan kesesuaian implementasi, tetapi tidak menjamin autentikasi integritas cipherteks. Percobaan manipulasi satu bit pada cipherteks menghasilkan respons dekripsi berstatus HTTP 200 tanpa deteksi kegagalan integritas, membuktikan bahwa penolakan galat saat ini hanya bertumpu pada pemeriksaan batas blok dan struktur padding. Implementasi belum menyertakan mekanisme autentikasi seperti \textit{Message Authentication Code} (MAC) atau skema \textit{Authenticated Encryption}.

  \item \textbf{Karakteristik Statistik dan Keterbatasannya:}
  \begin{itemize}
    \item \textit{Avalanche effect} satu blok menghasilkan rata-rata perubahan bit sebesar 49,976\% untuk pembalikan satu bit plainteks dan 49,557\% untuk pembalikan satu bit kunci (256 percobaan, seed 42), yang mendekati nilai ideal 50\% pada kriteria difusi empiris.
    \item Pada pesan utuh, perambatan perubahan bit plainteks mencerminkan sifat masing-masing mode operasi: ECB hanya mencapai 1,933\% karena pengaruh perubahan terbatas pada satu blok; CBC dan CFB mencapai 26,001\% dan 24,164\% akibat mekanisme \textit{chaining} dan umpan balik; sedangkan OFB dan CTR menghasilkan tepat satu bit perubahan (0,031\%) karena sifat \textit{keystream} independen. Sebaliknya, pembalikan bit kunci menghasilkan perubahan sekitar 50\% pada kelima mode (49,908\%--50,054\%) karena kunci memengaruhi seluruh pemanggilan cipher.
    \item Entropi empiris teks meningkat dari 4,860 bit/byte menjadi 7,188 bit/byte pada ECB dan 7,433--7,571 bit/byte pada mode lainnya. Pada berkas biner, mode non-ECB mencapai entropi 7,804--7,817 bit/byte dengan sebaran histogram yang mendekati seragam. Sebaliknya, ECB pada biner hanya mencapai 7,220 bit/byte dengan 86 nilai byte yang tidak muncul akibat blok masukan yang berulang.
    \item Kemerataan histogram dan tingginya entropi empiris merupakan indikator sebaran frekuensi byte pada sampel uji, tetapi bukan merupakan jaminan ketiadaan pola atau bukti ketahanan terhadap kriptanalisis formal.
  \end{itemize}
\end{enumerate}

\section{Saran Pengembangan}
Berdasarkan hasil evaluasi, keterbatasan pengujian, dan temuan implementasi saat ini, beberapa saran pengembangan yang diajukan adalah sebagai berikut:
\begin{enumerate}
  \item \textbf{Kajian Kriptanalisis Formal:} Melakukan analisis matematis mendalam terhadap ketahanan Brainrot Cipher terhadap kriptanalisis diferensial dan linier. Kajian lanjutan perlu menghitung probabilitas karakteristik diferensial maksimum (\textit{maximum differential characteristic probability}), menyusun tabel profil perbedaan (DDT) dan aproksimasi linier (LAT) Kotak-S $S_{67}$, serta membuktikan batas bawah jumlah putaran aktif (\textit{active S-boxes}) di bawah pengaruh kombinasi ShiftRows67, MixColumns67, Bit-Plane Transpose (BPT), dan langkah Feistel FCM.

  \item \textbf{Integrasi Autentikasi Integritas dan Format Berkas:}
  \begin{itemize}
    \item Mengembangkan mode operasi terotentikasi, seperti GCM (\textit{Galois/Counter Mode}) atau skema \textit{Encrypt-then-MAC} menggunakan HMAC atau KMAC, agar sistem dapat memverifikasi integritas dan keaslian cipherteks serta menolak manipulasi data secara kriptografis.
    \item Merancang format wadah berkas terstandardisasi yang menyertakan \textit{header} metadata (pengenal format, identitas mode, skema padding, serta nilai IV acak) sehingga proses dekripsi berkas tidak lagi memerlukan masukan parameter manual pengguna.
  \end{itemize}

  \item \textbf{Penyempurnaan Antarmuka dan Manajemen Kunci/IV:}
  \begin{itemize}
    \item Memperbaiki penanganan IV otomatis pada antarmuka web dan API agar nilai IV acak yang dibangkitkan oleh backend selalu dikembalikan kepada klien atau disematkan di awal berkas keluaran secara konsisten.
    \item Menghapus mekanisme \textit{fallback} simulasi XOR pada antarmuka peramban dan menggantinya dengan penanganan galat jaringan yang eksplisit guna mencegah kebingungan pengguna terhadap hasil enkripsi.
  \end{itemize}

  \item \textbf{Optimasi Kinerja dan Evaluasi Saluran Samping:} Menerapkan tabel pra-komputasi (\textit{T-tables}) yang menggabungkan operasi SubBytes67, ShiftRows67, dan MixColumns67 untuk meningkatkan \textit{throughput} pemrosesan perangkat lunak, serta mengevaluasi implementasi terhadap kerentanan saluran samping (\textit{side-channel attacks}), terutama serangan waktu (\textit{timing attack}) pada operasi perkalian medan berhingga $\mathrm{GF}(2^8)$ dan operasi percabangan bit.
\end{enumerate}

\clearpage\phantomsection
\addcontentsline{toc}{chapter}{DAFTAR PUSTAKA}
\begin{thebibliography}{99}
\bibitem{spesifikasi}
A. E. R. Xavier dan N. G. Subrata,
``Tugas 2 IF4020-Kriptografi Semester 1 Tahun 2026/2027:
Membuat Block Cipher Baru,'' spesifikasi tugas, 2026.
Salinan lokal: \path{Tugas2-IF4020-2026.md}.

\bibitem{aes}
National Institute of Standards and Technology,
``Advanced Encryption Standard (AES),'' FIPS 197-upd1,
9 Mei 2023. [Daring]. Tersedia:
\url{https://doi.org/10.6028/NIST.FIPS.197-upd1}.

\bibitem{mode}
M. Dworkin, ``Recommendation for Block Cipher Modes of Operation:
Methods and Techniques,'' NIST SP 800-38A, Desember 2001.
[Daring]. Tersedia: \url{https://doi.org/10.6028/NIST.SP.800-38A}.

\bibitem{padding}
R. Housley, ``Cryptographic Message Syntax (CMS),'' RFC 5652,
bagian 6.3 dan 9, September 2009. [Daring]. Tersedia:
\url{https://www.rfc-editor.org/rfc/rfc5652}.

\bibitem{shannon}
C. E. Shannon, ``Communication Theory of Secrecy Systems,''
\textit{Bell System Technical Journal}, vol. 28, no. 4,
hlm. 656--715, Oktober 1949. [Daring]. Tersedia:
\url{https://doi.org/10.1002/j.1538-7305.1949.tb00928.x}.

\bibitem{heys}
H. M. Heys, ``A Tutorial on Linear and Differential Cryptanalysis,''
laporan teknis CORR 2001-17, Centre for Applied Cryptographic Research,
University of Waterloo, Maret 2001. [Daring]. Tersedia:
\url{https://www.engr.mun.ca/~howard/PAPERS/ldc_tutorial.pdf}.

\bibitem{cryptosys}
D. Ireland, ``Padding schemes for block ciphers,''
\textit{CryptoSys PKI Pro Manual}, DI Management Services,
2026. [Daring]. Tersedia:
\url{https://www.cryptosys.net/pki/manpki/pki_paddingschemes.html}.
\end{thebibliography}

% Tambahkan sumber yang benar-benar dibaca/digunakan; kutip dengan \cite{...}.

\BabTanpaNomor{LAMPIRAN}
\newcounter{bagianlampiran}
\newcommand{\BagianLampiran}[1]{%
  \refstepcounter{bagianlampiran}
  \section*{\Alph{bagianlampiran}. #1}\phantomsection
  \addcontentsline{toc}{section}{\Alph{bagianlampiran}. #1}%
}

\BagianLampiran{Pranala Repositori}
\expandafter\url\expandafter{\Repositori}

\Isi{Cantumkan tag release atau commit yang sesuai dengan hasil laporan.}

\BagianLampiran{Video Demo dan Dokumentasi API (Opsional)}
Video demo: \url{https://youtu.be/k17aw6r5oJw}.

Dokumentasi API menggunakan Swagger UI pada
\url{http://localhost:8000/docs/}, dengan spesifikasi OpenAPI pada
\url{http://localhost:8000/openapi.json}. Parameter, keluaran, kesalahan,
dan contoh permintaan teks/berkas tersedia pada \path{docs/api.md}
dan README repositori. Dokumentasi mencakup health check, pemrosesan
cipher, trace, dan avalanche.
\Isi{Tambahkan pranala Swagger UI publik setelah dokumentasi dipublikasikan.}

\BagianLampiran{Pembagian Tugas Anggota}
\begin{table}[H]
  \centering\small
  \caption{Isian pembagian tugas kelompok.}
  \begin{tabularx}{\linewidth}{@{}p{0.32\linewidth}p{0.18\linewidth}Y@{}}
    \toprule Nama & NIM & Kontribusi\\ \midrule
    \Muhammad Ra'if Alkaustar & \13523011 & \Isi{Perancangan arsitektur inti cipher, ekspansi subkunci (key schedule), pengujian performa benchmark, penulisan laporan analisis keamanan LaTeX.}\\[4pt]
    \Muhammad Aditaya Rahmadeni & \13523028 & \Isi{Pengembangan antarmuka Web Neo-Brutalist (React + TypeScript + Tailwind), visualisasi metrik analitik (Avalanche, Entropi, Histogram), arsitektur REST API & routing HTTP, kontainerisasi Docker & CI/CD deployment}\\[4pt]
    \Fityatul Haq Rosyidi & \13523116 & \Isi{Implementasi transformasi cipher (generasi tabel S-Box dinamis & ShiftRows), utilitas pengujian Command Line Interface (CLI Crypt), integrasi layanan enkripsi/dekripsi, penyusunan dokumentasi dasar teori.}\\[4pt]
    \bottomrule
  \end{tabularx}
\end{table}

\BagianLampiran{Petunjuk Menjalankan}
Jalankan dari root repositori dengan Go 1.23 atau lebih baru.
\begin{lstlisting}[language={},numbers=none,caption={Perintah verifikasi dan analisis.}]
go version
go -C src test ./...
go -C src run ./cmd/crypt -op encrypt -mode ECB \
  -key-hex 00112233445566778899aabbccddeeff -input "Halo Brainrot"
go -C src run ./cmd/analyze -cipher brainrot -modes all \
  -trials 256 -seed 42 -output ../reports/brainrot.json
\end{lstlisting}
Di PowerShell, satukan setiap perintah menjadi satu baris atau gunakan
backtick sebagai penyambung baris. README repositori memuat contoh
enkripsi/dekripsi dengan variabel PowerShell.

Backend web dijalankan dengan \texttt{go -C src run ./cmd/server},
sedangkan frontend memakai \texttt{npm --prefix frontend run dev}
setelah dependensinya dipasang. Demonstrasi pemulihan data dan kasus
perubahan dapat diulang dengan
\texttt{go -C src run ../scripts/report\_demonstration.go}.
Pengelolaan IV dan batas verifikasi integritas dijelaskan pada Bab IV.

\end{document}
